\chapter{Discussion}
\label{chap:discussion}

\section{Interpretation}
\label{sec:disc-interp}
The central finding of this thesis is that a compact VLA can be fine-tuned and
deployed within a 12~GB GPU budget to perform a language-conditioned manipulation
task with an 80\% task-completion rate. Two aspects of \emph{how} that result was
obtained are as important as the number itself.

First, the largest improvements came from \emph{methodology, not model scale}. The
preliminary study (\refsec{sec:results-journey}) showed that fixing a success-metric
bug and removing duplicated demonstrations moved real performance far more than any
architectural change, and the supermarket result was reached with a
\emph{frozen} backbone and only 208 training demonstrations. This is consistent with the
literature: $\pi_{0.5}$ and RoboAgent both attribute performance to data diversity
and composition rather than to scaling the action model alone
\parencite{black2025pi05, bharadhwaj2024roboagent}. The frozen-backbone strategy
also validates the pretrain-then-post-train view \parencite{xiang2026posttraining}:
the model arrives already knowing what the objects are and what the words mean, and
fine-tuning need only supply the motor mapping --- which is why so few
demonstrations sufficed.

Second, the 0\%-to-80\% improvement (\refsec{sec:results-diagnosis}) came from
\emph{diagnosis rather than brute force}. Tracing a single failure revealed that
the policy already understood the task and was near-missing at the final drop; the
fixes targeted observability and episode boundaries specifically. This is a more
transferable outcome than a number obtained by simply retraining a larger model.

The language-grounding test (\refsec{sec:results-language}) confirms that the
policy is genuinely instruction-following: with the scene fixed, changing only the
instruction word moves the reach to the named item. Together with the graded
per-item success rates, this rules out the hypothesis that the policy executes a
\emph{single} memorised trajectory irrespective of the instruction.

It does not, however, rule out the hypothesis that the policy executes \emph{one
memorised trajectory per item word}, and \refsec{sec:results-lookup} shows that this
is what it does. The distinction is the most consequential finding of this thesis, and
it is worth being precise about why the weaker test was insufficient. Every item
occupied a fixed position in all 208 training episodes and in the grounding test
itself; under those conditions ``localise the named object and reach to it'' and
``recall where that word's object sits and reach there'' are behaviourally identical.
Only moving the objects separates them, and when they are moved the policy follows the
recalled position without exception. The lesson generalises beyond this thesis: a
grounding test conducted on a positionally static scene cannot establish grounding, no
matter how cleanly the reach tracks the instruction, and results of that form in the
wider literature should be read with the same reservation. That this thesis's own
earlier claim did not survive the stronger test is an argument for running it.

The success of the frozen-backbone strategy deserves emphasis, because it reframes
what fine-tuning a VLA on a small dataset actually involves. The intuition that a
policy must ``learn to see'' the objects is misleading in this setting: the backbone
already sees and reads competently from its pretraining, and the 208 training demonstrations
supply only the missing ingredient --- how this particular embodiment should move to
act on what is already perceived. This division of labour is what makes the data
requirement so modest, and it explains why an attempt to retrain the perception on so
few examples was not merely unnecessary but actively harmful in the pilot, where it
either exhausted memory or destabilised training. The broader lesson is that, under
constraint, the right question is not ``how much of the model can I train?'' but ``how
little of it \emph{must} I train?''. Identifying the smallest sufficient set of
parameters to adapt --- here, the action expert alone --- is both a memory strategy and a
regularisation strategy, and it is a transferable principle for fine-tuning large
pretrained models on limited data in general, not only for this task.

A final interpretive point concerns the decomposition of the task into grasping and
placing. The results consistently show grasping to be the more nearly solved
sub-problem (92.5\% overall, 100\% on three of four items) and placing to be the more
variable one (80\% overall). This is initially counter-intuitive, since grasping ---
aligning a gripper with an object and securing it --- is usually considered the harder
manipulation primitive. The explanation is specific to this task: the grasp is
executed close to the item with a clear eye-in-hand view and benefits directly from
the scripted expert's carefully tuned per-item poses, whereas the placement is a
release into a small target at the end of a longer, stochastic motion, where residual
error has accumulated and (before the basket camera) observability was poorest. The
implication is that, for this class of task, the marginal engineering effort is better
spent on the terminal placement --- its observability, its target tolerance, and the
data around it --- than on the grasp, and this is precisely where the 0\%-to-80\% fixes
were directed. It is a small but useful design heuristic that generalises beyond the
present system: instrument and ease the least-observed, highest-variance phase of the
motion first.

\section{Limitations}
\label{sec:disc-limits}
This work has clear and honestly stated limitations.

\subsection{Weak item.} The milk carton is the weakest item (70\% grasp, 60\%
placement). It is the tallest item, so its peak lift clears the grasp threshold by the smallest margin of any item, and the learned reach exhibits a consistent inward bias (Sec. 4.3). It is not the most laterally extended target — the water bottle sits 8 cm further out and grasps at 100 \%. It placed more reliably at earlier checkpoints, so the weakness is recoverable with milk-specific grasp tuning or additional milk demonstrations rather than being fundamental.

\subsection{Multi-item lists.} Single-item picks are reliable, but whole-list
completion is not: 45\% for two items and 30\% for three
(\refsec{sec:results-orchestrator}). An accumulating basket was previously assumed to
be the cause, on the reasoning that the policy never saw a non-empty basket in
training. Measurement does not support this: success does not decline as the basket
fills. The mechanism remains open, and the most plausible candidate is the one
identified in \refsec{sec:results-lookup} --- earlier picks nudge the remaining items,
and a policy that reaches a recalled position rather than a perceived one fails once
its target has moved.

\subsection{Reliance on fixed object positions.} The policy does not localise the named
item; it recalls where that item was demonstrated (\refsec{sec:results-lookup}). This
is the root cause behind several entries in this list, including the \SI{2.5}{cm}
robustness result and, most likely, the multi-item degradation above. It is a data
limitation rather than a model limitation: with every item at a fixed position in all
208 episodes, a word-to-trajectory association fits the demonstrations exactly and is
far simpler to learn than visual search, so the training set never gave the policy a
reason to look. The direct remedy --- re-collecting demonstrations with item positions
randomised over a substantial range, so that the recalled-position strategy no longer
fits the data --- is the single highest-value next experiment this work suggests, and is
developed in \refsec{sec:disc-ideal}.

\subsection{Excluded item.} The cereal box is excluded from training because its
wide face slips in the two-finger gripper during transport; a wider grasp or a
different grip strategy would be required.

\subsection{Simulation only.} All results are in simulation; no sim-to-real transfer
is claimed. The pipeline is deliberately real-robot-shaped, however --- it uses the
LeRobot framework and accurate hardware models (UR10e, Robotiq) --- so the same
code path could fine-tune on real demonstrations.

\subsection{Stochastic policy.} The flow-matching action expert samples actions
stochastically, so there is run-to-run variance, especially on the weak items. This
is exactly why success is reported over 320 trials with confidence intervals rather
than from single runs.

\subsection{Checkpoint selection.} Checkpoints are selected by closed-loop success on
evaluation episodes that use fresh seeds and paraphrases but the same item set; a
fully locked, never-used-for-selection test set would be a cleaner final estimate.

\subsection{Inference cost and control rate.} Although the frozen-backbone strategy
makes \emph{training} inexpensive, \emph{inference} still runs the full
vision--language backbone every re-planning step, which bounds the achievable control
rate on modest hardware. This is acceptable for the quasi-static manipulation task
studied here, where the environment does not change during a pick, but it would be a
first-order concern for the dynamic navigation stage, where a slow control loop is
unsafe around moving people. This is a further, concrete argument for the dual-system
architecture discussed in \refsec{sec:disc-ideal}, in which only a small reflex policy
need run at high rate. The action-chunk horizon is a further lever: executing more of
each predicted chunk before re-planning amortises the backbone cost but reduces
reactivity, the non-monotonic trade-off noted in \refsec{sec:bg-eval}.

\subsection{Scope of the instruction set.} The language conditioning is exercised over
six paraphrase templates spanning four items --- 24 distinct instructions; it is not
tested on open-vocabulary requests, unseen object names, compositional instructions
(``the small red can''), or referring expressions that require disambiguation. The
grounding result of \refsec{sec:results-language} should therefore be read as evidence
that the policy uses the instruction to select among four known items, and
\refsec{sec:results-lookup} narrows this further: the selection is among four
memorised behaviours, not among four perceived objects. Neither result supports a
claim of open-world language understanding, which would require the far more diverse
data discussed as future work.

\section{On cross-benchmark evaluation and specialisation}
\label{sec:disc-crosseval}
A natural question is whether the supermarket-fine-tuned policy could be evaluated
directly on the RoboCasa kitchen benchmark, to compare it against the reference
foundation model. It could be run, but the result would not be meaningful, because
fine-tuning specialises the policy to a specific \emph{interface} as well as a
specific task. The action expert and its frozen normalisation statistics are tied to
this project's embodiment: a seven-dimensional action of six UR10e joint targets
plus a gripper command, and three cameras (wrist, scene, basket) at a fixed
resolution and placement. RoboCasa uses a different embodiment (a Franka arm on a
mobile base), a different action convention and dimensionality, and different camera
streams and objects. Feeding RoboCasa observations to this policy and applying its
outputs to a Franka would be an interface mismatch, not a capability test; failure
there would say nothing about the policy's competence, only that it was asked to
speak a language it was never trained in. This is why the meaningful RoboCasa
experiment reported in \refsec{sec:results-journey} is the converse one --- fine-tune
the small models \emph{on} RoboCasa demonstrations and evaluate \emph{on} RoboCasa,
which yielded 0/20 --- and why generalisation of the present policy is better assessed
on held-out variations of its own task (unseen positions, novel paraphrases) than on
a foreign benchmark. Cross-embodiment transfer is a distinct research problem,
addressed in the literature by co-training across embodiments
\parencite{black2025pi05}, and is out of scope here.

\section{Toward a real-world system without constraints}
\label{sec:disc-ideal}
It is worth asking what an ideal system would look like if compute were not a
constraint. Notably, the answer is not simply ``a larger model'': some limits are
physical rather than computational. A very large model cannot run a fast control
loop, which is unsafe around moving people. The field therefore converges on a
\emph{dual-system} design --- a large reasoner (``System~2'') for planning the list
and interpreting a dynamic scene, paired with a small, fast reflex policy
(``System~1'') for real-time control \parencite{nvidia2025gr00t, black2025pi05}.

Beyond the architecture, five changes would most improve a real deployment; each is
motivated directly by a limitation or observation of the present work. The first is
both the cheapest and, on the evidence of \refChap{chap:results}, the most valuable.

\subsection{Positional randomisation in the demonstrations.} The single most
consequential change would require no additional hardware, no larger model and no new
algorithm: re-collect the demonstrations with each item's position randomised over a
substantial range rather than fixed. \refsec{sec:results-lookup} shows that the policy
currently maps each item word to a memorised trajectory, and
\refsec{sec:results-robustness} shows the cost of that --- a \SI{2.5}{cm} displacement
removes more than twenty points of task completion. Neither is a failing of the model:
with every item at a fixed position in all 208 episodes, recall fits the training data
exactly, so the dataset never made visual localisation necessary. Randomising positions
over, say, the full width of the shelf would make the recalled-position strategy
inconsistent with the demonstrations and force the policy to use the visual information
it already receives --- information the scripted expert demonstrably uses, since it
re-solves its grasp inverse kinematics on the settled object position and maintains
96--100\% per-item success under exactly the jitter that defeats the learned policy.
The experiment is inexpensive, since the data-collection pipeline already supports
per-episode jitter and need only have its magnitude increased, and it is the natural
next step for this work. Its outcome would also be informative either way: if
randomisation produces genuine localisation, the limitation is confirmed as one of
data; if it does not, the limitation lies deeper, in the capacity of a frozen backbone
with a small action expert to support visual search.

\subsection{Whole-body control.} Controlling the base and arm as a single system,
rather than parking the base and then moving the arm, would let the robot reach while
repositioning --- extending its effective workspace, reaching around obstacles, and
approaching items from a more favourable base pose. The parked-base assumption made
here simplifies the manipulation problem: every item must be reached from a single
fixed stance, so the policy has one approach geometry per item and no means of
improving it. A whole-body policy would fold base positioning into the same
optimisation as the reach, allowing a poor approach angle to be corrected by moving
rather than by stretching.

\subsection{Tactile and force sensing.} The clearest capability gap in this thesis is
grasping: the milk carton grasps at 70\% and the cereal box is excluded entirely for
slipping. These are difficult to solve from vision alone, because whether the fingers
have secured an object --- and whether it is slipping during transport --- is a contact
property, not a visual one. Tactile or force sensing on the gripper would provide
exactly this signal, enabling closed-loop grip adjustment and slip detection, and is
the most direct route to lifting the weak items and re-admitting cereal.

\subsection{Diverse, co-trained data.} A four-item, single-scene dataset cannot
support the generalisation a service robot requires. The most effective route,
following $\pi_{0.5}$ \parencite{black2025pi05}, is to co-train on a deliberately
heterogeneous mixture --- teleoperation across many stores and items, domain-randomised
simulation, and even human video --- so that the policy learns invariances rather than
memorising a fixed inventory. It would also broaden the scene distribution generally, though note that basket occupancy was measured not to be the cause of the multi-item degradation (Sec. 4.11).

\subsection{Imitation followed by reinforcement.} Behaviour cloning caps a policy at
the quality of its demonstrations; it can, at best, imitate the scripted expert.
Reinforcement-learning fine-tuning on top of the imitation-pretrained policy ---
safely, in simulation --- could discover behaviours the expert never demonstrated, in
particular robust recoveries from the near-misses that dominate the residual failures
here, pushing performance past the demonstration ceiling.

In this framing, the system built in this thesis is the \emph{deployable subset} of
that ideal: the same fundamental recipe --- a pretrained perception prior, a learned
action mapping, and disciplined evaluation --- constrained to what a single GPU and a
modest demonstration budget allow. The value of characterising that subset carefully
is that it establishes what is achievable now, and identifies precisely which
additions --- contact sensing, whole-body control, data diversity, reinforcement ---
would buy the most, and why.

\section{Comparison with related work}
\label{sec:disc-compare}
The 80\% task-completion rate sits in the same regime as Kiviranta's 80\% single-task
success fine-tuning the same SmolVLA checkpoint on real SO-101 hardware from around
100 demonstrations \parencite{kiviranta2025smolvla}. The two figures should not be
equated, however: Kiviranta's is obtained on physical hardware, where object pose,
lighting and contact vary from trial to trial, whereas the figure here is measured in
simulation at a fixed object pose, and falls to 58.8\% once object position is allowed
to vary by \SI{2.5}{cm} (\refsec{sec:results-robustness}). The coincidence of the two
numbers is therefore a coincidence of regime, not of difficulty. The comparison is
nonetheless consistent with that work's caution that SmolVLA's headline numbers are
easy to overstate without careful evaluation --- a caution this thesis addresses
directly through Wilson intervals, adequate sample sizes, success-based checkpoint
selection and, ultimately, by testing and revising its own headline claim. The observation that
data and metrics dominate model choice at this scale aligns with
\textcite{black2025pi05} and \textcite{bharadhwaj2024roboagent}, and the sensitivity
of chunked policies to execution timing observed in the failure analysis is
consistent with \textcite{wang2026vlalimits}.

It is worth being precise about what this thesis adds relative to that body of work.
It does not propose a new model, a new architecture, or a new training algorithm;
SmolVLA, flow matching and behaviour cloning are all taken from the literature. Its
contribution is instead an \emph{integrated, rigorously evaluated case study} under an
explicit and unusually tight resource constraint, of a kind that is comparatively rare
in a field that often reports headline capabilities from large systems. The value of
such a study is threefold: it establishes a concrete, reproducible baseline for what a
compact VLA achieves on a controlled language-conditioned task within 12~GB; it
demonstrates a methodology --- honest metrics, adequate samples, confidence intervals,
success-based selection, and failure diagnosis --- that is transferable to other
constrained settings; and it produces a first-hand data point on the model-versus-data
question that corroborates, from an independent direction, the conclusions of larger
studies \parencite{black2025pi05, bharadhwaj2024roboagent}. In a research culture that
can over-weight scale, a careful account of what is achievable \emph{without} it is a
useful counterweight.

\section{Threats to validity}
\label{sec:disc-validity}
It is worth stating explicitly the threats to the validity of the conclusions and
how each is mitigated or bounded.

\subsection{Construct validity.} Task completion is defined as the item coming to
rest inside the basket footprint and low, rather than merely near the basket. This
was a deliberate correction of the proxy metric that inflated results in the pilot
study (\refsec{sec:results-journey}), and it ensures the headline number measures
the intended construct --- successful placement --- rather than a correlate of it.

\subsection{Internal validity.} The 0\%-to-80\% improvement is attributed to three
specific changes made together (clean episode boundaries, a basket camera and a
wider tote), so their individual contributions are not isolated by an ablation. The
changes were, however, motivated by a concrete traced failure, and each targets a
distinct facet of that failure (what the policy is taught, what it can see, and how
much slack the target allows). A controlled ablation is identified as useful further
work.

\subsection{External validity.} All results are in a single simulated environment
with four items and a parked base; they should not be read as a claim about real
hardware, other objects, or full mobile manipulation. A further and more specific
bound applies: the headline evaluation places each item at a fixed pose, so it
measures performance on the exact configuration demonstrated rather than on a
distribution of configurations. \refsec{sec:results-robustness} quantifies the
difference, and \refsec{sec:results-lookup} explains it. Any generalisation from these
results to a setting in which objects are not where the policy last saw them is
unsupported. The intent is narrower --- to characterise what a compact VLA can do under
a fixed hardware budget on a controlled task --- and the limitations of
\refsec{sec:disc-limits} bound the generalisation accordingly.

\subsection{Statistical conclusion validity.} Because the flow-matching policy is
stochastic, single rollouts are uninformative; success is therefore reported over
320 trials with Wilson intervals, and differences whose intervals overlap (for
example the 5k and 10k checkpoints) are not treated as real. The intervals are wide
at $n=20$ per item, so per-item claims (especially milk at 60\%) are made
cautiously. A limitation of the checkpoint-selection procedure is that selection and
reporting share the same item set; a fully locked test set would give an unbiased
final estimate.
